\chapter{2013年广东卷（文）}

\section{单选题}

\begin{problem}
设集合 \(S = \{x \mid x^2 + 2x = 0, x \in \R\}\)，\(T = \{x \mid x^2 - 2x = 0, x \in \R\}\)，则 \(S \cap T =\)（\quad）

\choices
    {\( \{0\} \)}
    {\(\{0, 2\}\)}
    {\(\{-2, 0\}\)}
    {\(\{-2, 0, 2\}\)}
\end{problem}

\begin{answer}
A
\end{answer}

\begin{solution}
由 \(x^2+2x=x(x+2)=0\)，得 \(S=\{-2,0\}\)；由 \(x^2-2x=x(x-2)=0\)，得 \(T=\{0,2\}\). 因此 \(S\cap T=\{0\}\)，故选 A.
\end{solution}

\begin{problem}
函数 \( y = \frac{\lg(x + 1)}{x - 1} \) 的定义域是（\quad）

\choices
    {\( (-1,+\infty) \)}
    {\( [-1,+\infty) \)}
    {\( (-1,1)\cup(1,+\infty) \)}
    {\([-1, 1) \cup (1, +\infty)\)}
\end{problem}

\begin{answer}
C
\end{answer}

\begin{solution}
对数的真数必须为正，因此 \(x+1>0\)，即 \(x>-1\)；分式的分母不能为零，因此 \(x\neq1\). 所以定义域为 \((-1,1)\cup(1,+\infty)\)，故选 C.
\end{solution}

\begin{problem}
若 \(\i (x + y\i) = 3 + 4\i\)，\(x\)，\(y \in \R\)，则复数 \( x + y\i \) 的模是（\quad）

\choices
    {\(2\)}
    {\(3\)}
    {\(4\)}
    {\(5\)}
\end{problem}

\begin{answer}
D
\end{answer}

\begin{solution}
因为 \(\i^2=-1\)，所以 \(\i(x+y\i)=-y+x\i\). 与 \(3+4\i\) 比较实部和虚部，得 \(x=4\)，\(y=-3\). 于是 \(|x+y\i|=|4-3\i|=\sqrt{4^2+(-3)^2}=5\)，故选 D.
\end{solution}

\begin{problem}
已知 \(\sin \left(\frac{5\pi}{2} + \alpha\right) = \frac{1}{5}\)，那么 \(\cos \alpha =\)（\quad）

\choices
    {\(-\frac{2}{5}\)}
    {\(-\frac{1}{5}\)}
    {\(\frac{1}{5}\)}
    {\(\frac{2}{5}\)}
\end{problem}

\begin{answer}
C
\end{answer}

\begin{solution}
由诱导公式，\(\sin\left(\frac{5\pi}{2}+\alpha\right)=\sin\left(\frac{\pi}{2}+\alpha\right)=\cos\alpha\). 题设给出该值为 \(\frac15\)，所以 \(\cos\alpha=\frac15\)，故选 C.
\end{solution}

\begin{problem}
执行如图所示的程序框图，若输入 \(n\) 的值为 3 ，则输出 \(s\) 的值是（\quad）

\bitmapfigure[max width=0.34\linewidth]{img/2013/guangdong_liberal/q5_fig1.png}

\choices
    {\(1\)}
    {\(2\)}
    {\(4\)}
    {\(7\)}
\end{problem}

\begin{answer}
C
\end{answer}

\begin{solution}
初始时 \(i=1\)，\(s=1\)，且循环条件为 \(i\leq n\). 当 \(n=3\) 时，第一次循环后 \(s=1+(1-1)=1\)，\(i=2\)；第二次循环后 \(s=1+(2-1)=2\)，\(i=3\)；第三次循环后 \(s=2+(3-1)=4\)，\(i=4\). 此时 \(4\leq3\) 不成立，输出 \(s=4\)，故选 C.
\end{solution}

\begin{problem}
某三棱锥的三视图如图所示，则该三棱锥的体积是（\quad）

\bitmapfigure[max width=0.48\linewidth]{img/2013/guangdong_liberal/q6_fig1.png}

\choices
    {\( \frac{1}{6} \)}
    {\( \frac{1}{3} \)}
    {\( \frac{2}{3} \)}
    {\(1\)}
\end{problem}

\begin{answer}
B
\end{answer}

\begin{solution}
由三视图可知，该三棱锥的高为 \(2\)，底面为等腰直角三角形，两个直角边长均为 \(1\). 因此体积为
\[
V=\frac13\times\left(\frac12\times1\times1\right)\times2=\frac13
\]
故选 B.
\end{solution}

\begin{problem}
垂直于直线 \( y = x + 1 \) 且与圆 \( x^{2} + y^{2} = 1 \) 相切于第一象限的直线方程是（\quad）

\choices
    {\(x + y - \sqrt{2} = 0\)}
    {\(x + y + 1 = 0\)}
    {\(x + y - 1 = 0\)}
    {\(x + y + \sqrt{2} = 0\)}
\end{problem}

\begin{answer}
A
\end{answer}

\begin{solution}
直线 \(y=x+1\) 的斜率为 \(1\)，所求直线垂直于它，故斜率为 \(-1\)，可设方程为 \(x+y+c=0\). 它与单位圆相切，所以原点到该直线的距离为 \(1\)，即
\[
\frac{|c|}{\sqrt2}=1
\]
从而 \(c=\pm\sqrt2\). 第一象限的切点为 \(\left(\frac1{\sqrt2},\frac1{\sqrt2}\right)\)，代入 \(x+y+c=0\) 得 \(c=-\sqrt2\). 所以所求直线为 \(x+y-\sqrt2=0\)，故选 A.
\end{solution}

\begin{problem}
设 \(l\) 为直线，\(\alpha\)，\(\beta\) 是两个不同的平面，下列命题中正确的是（\quad）

\choices
    {若 \(l \parallel \alpha, l \parallel \beta\)，则 \(\alpha \parallel \beta\)}
    {若 \(l \perp \alpha\)，\(l \perp \beta\)，则 \(\alpha \parallel \beta\)}
    {若 \(l \perp \alpha\)，\(l \parallel \beta\)，则 \(\alpha \parallel \beta\)}
    {若 \(\alpha \perp \beta\)，\(l \parallel \alpha\)，则 \(l \perp \beta\)}
\end{problem}

\begin{answer}
B
\end{answer}

\begin{solution}
命题 A 不成立：取 \(\alpha\colon z=0\)、\(\beta\colon x=0\) 和直线 \(l\colon (x,y,z)=(1,t,1)\). 直线 \(l\) 的方向平行于两平面的交线方向，因而 \(l\parallel\alpha\) 且 \(l\parallel\beta\)，但 \(\alpha\) 与 \(\beta\) 相交而不平行.

命题 B 成立：若一条直线同时垂直于两个平面，则它同时是两个平面的法线方向，从而两个平面的法向量平行. 由于两平面不同，所以两平面平行.

命题 C 不成立：取 \(\alpha\colon z=0\)，\(\beta\colon y=0\)，直线 \(l\colon (x,y,z)=(1,1,t)\). 此时 \(l\perp\alpha\)，且 \(l\parallel\beta\)，但 \(\alpha\) 与 \(\beta\) 垂直而不平行. 命题 D 也不成立：取 \(\alpha\colon z=0\)，\(\beta\colon x=0\)，直线 \(l\colon (x,y,z)=(1,t,1)\). 此时 \(\alpha\perp\beta\)，且 \(l\parallel\alpha\)，但 \(l\) 的方向向量平行于 \(\beta\)，所以 \(l\) 不垂直于 \(\beta\). 因此正确命题只有 B.
\end{solution}

\begin{problem}
已知中心在原点的椭圆 \(C\) 的右焦点为 \(F(1,0)\). 离心率等于 \(\frac{1}{2}\)，则 \(C\) 的方程是（\quad）

\choices
    {\(\frac{x^2}{3} +\frac{y^2}{4} = 1\)}
    {\(\frac{x^2}{4} +\frac{y^2}{\sqrt{3}} = 1\)}
    {\(\frac{x^2}{4} +\frac{y^2}{2} = 1\)}
    {\(\frac{x^2}{4} +\frac{y^2}{3} = 1\)}
\end{problem}

\begin{answer}
D
\end{answer}

\begin{solution}
椭圆中心在原点且焦点为 \(F(1,0)\)，故焦距参数 \(c=1\). 由离心率 \(e=\frac ca=\frac12\)，得长半轴 \(a=2\). 于是
\[
b^2=a^2-c^2=4-1=3
\]
椭圆的焦点在 \(x\) 轴上，所以其方程为 \(\frac{x^2}{4}+\frac{y^2}{3}=1\)，故选 D.
\end{solution}

\begin{problem}
设 \(\bs{a}\) 是已知的平面向量且 \(\bs{a} \neq 0\). 关于向量 \(\bs{a}\) 的分解，有如下四个命题：
\begin{circlelist}
\item 给定向量 \(\bs{b}\)，总存在向量 \(\bs{c}\)，使 \(\bs{a} = \bs{b} + \bs{c}\).
\item 给定向量 \(\bs{b}\) 和 \(\bs{c}\)，总存在实数 \(\lambda\) 和 \(\mu\)，使 \(\bs{a} = \lambda \bs{b} + \mu \bs{c}\)；
\item 给定单位向量 \(\bs{b}\) 和正数 \(\mu\)，总存在单位向量 \(\bs{c}\) 和实数 \(\lambda\)，使 \(\bs{a} = \lambda \bs{b} + \mu \bs{c}\).
\item 给定正数 \(\lambda\) 和 \(\mu\)，总存在单位向量 \(\bs{b}\) 和单位向量 \(\bs{c}\)，使 \(\bs{a} = \lambda \bs{b} + \mu \bs{c}\).
\end{circlelist}

上述命题中的向量 \(\bs{b}\)，\(\bs{c}\) 和 \(\bs{a}\) 在同一平面内且两两不共线，
则真命题的个数是（\quad）

\choices
    {\(1\)}
    {\(2\)}
    {\(3\)}
    {\(4\)}
\end{problem}

\begin{answer}
B
\end{answer}

\begin{solution}
命题 1 中，对任意给定的 \(\bs{b}\)，取 \(\bs{c}=\bs{a}-\bs{b}\)，便有 \(\bs{a}=\bs{b}+\bs{c}\)，所以命题 1 为真.

由于 \(\bs{b}\) 与 \(\bs{c}\) 不共线，它们是同一平面的一组基，因此同一平面内的向量 \(\bs{a}\) 可唯一表示为 \(\bs{a}=\lambda\bs{b}+\mu\bs{c}\)，命题 2 为真.

命题 3 不总成立. 例如取 \(\bs{a}=(0,2)\)，单位向量 \(\bs{b}=(1,0)\)，并取 \(\mu=1\). 若存在单位向量 \(\bs{c}\) 和实数 \(\lambda\) 使 \(\bs{a}=\lambda\bs{b}+\bs{c}\)，则 \(\bs{c}=(-\lambda,2)\)，从而 \(|\bs{c}|=\sqrt{\lambda^2+4}\geq2\)，不可能等于 \(1\). 所以命题 3 为假.

命题 4 也不总成立. 对给定的非零向量 \(\bs{a}\)，取 \(\lambda=\mu=\frac14|\bs{a}|\). 若 \(\bs{b}\)、\(\bs{c}\) 都是单位向量，则由三角不等式有 \(|\lambda\bs{b}+\mu\bs{c}|\leq\lambda+\mu=\frac12|\bs{a}|<|\bs{a}|\)，不可能等于 \(\bs{a}\). 所以命题 4 为假. 真命题有 \(2\) 个，故选 B.
\end{solution}

\section{填空题}

\begin{problem}
设数列 \( \{a_{n}\} \) 是首项为 1，公比为 \(-2\) 的等比数列，则 \( a_{1} + |a_{2}| + a_{3} + |a_{4}| = \)\fillinblank{}.
\end{problem}

\begin{answer}
\(15\)
\end{answer}

\begin{solution}
等比数列的通项为 \(a_n=1\times(-2)^{n-1}\)，所以 \(a_1=1\)，\(a_2=-2\)，\(a_3=4\)，\(a_4=-8\). 因此
\[
a_1+|a_2|+a_3+|a_4|=1+2+4+8=15
\]
\end{solution}

\begin{problem}
若曲线 \( y = ax^2 - \ln x \) 在点 \((1, a)\) 处的切线平行于 \( x \) 轴，则 \( a = \)\fillinblank{}.
\end{problem}

\begin{answer}
\(\frac{1}{2}\)
\end{answer}

\begin{solution}
函数的定义域为 \(x>0\)，且 \(y'=2ax-\frac1x\). 在点 \((1,a)\) 处切线平行于 \(x\) 轴，故切线斜率为零，即
\[
y'(1)=2a-1=0
\]
解得 \(a=\frac12\).
\end{solution}

\begin{problem}
已知变量 \(x\)，\(y\) 满足约束条件 \(\begin{cases}
x - y + 3 \geqslant 0,\\
-1 \leqslant x \leqslant 1,\\
y \geqslant 1,
\end{cases}\) 则 \(z = x + y\) 的最大值是\fillinblank{}.
\end{problem}

\begin{answer}
\(5\)
\end{answer}

\begin{solution}
由 \(x-y+3\geq0\) 得 \(y\leq x+3\). 因此
\[
z=x+y\leq x+(x+3)=2x+3\leq5
\]
其中最后一个不等式利用了 \(x\leq1\). 当 \(x=1\)，\(y=4\) 时，三个约束均满足，且 \(z=1+4=5\)，所以最大值为 \(5\).
\end{solution}

\section{选考题：填空题}

\begin{problem}
已知曲线 \(C\) 的极坐标方程为 \(\rho = 2\cos \theta\)，以极点为原点，极轴为 \(x\) 轴的正半轴建立直角坐标系，则曲线 \(C\) 的参数方程为\fillinblank{}.
\end{problem}

\begin{answer}
\(
\begin{cases}
x=1+\cos\alpha,\\
y=\sin\alpha
\end{cases}
\)（\(\alpha\) 为参数）
\end{answer}

\begin{solution}
由极坐标关系 \(x=\rho\cos\theta\)，\(y=\rho\sin\theta\)，将 \(\rho=2\cos\theta\) 两边乘以 \(\rho\)，得
\[
\rho^2=2\rho\cos\theta=2x
\]
因为 \(\rho^2=x^2+y^2\)，所以曲线满足 \(x^2+y^2=2x\)，即
\[
(x-1)^2+y^2=1
\]
这是圆心为 \((1,0)\)、半径为 \(1\) 的圆，故可取参数方程
\[
\begin{cases}
x=1+\cos\alpha,\\
y=\sin\alpha,
\end{cases}
\]
其中 \(\alpha\) 为参数.
\end{solution}

\begin{problem}
如图，在矩形 \(ABCD\) 中，\(AB = \sqrt{3}\)，\(BC = 3\)，\(BE \perp AC\)，垂足为 \(E\)，则 \(ED = \)\fillinblank{}.

\bitmapfigure[max width=0.38\linewidth]{img/2013/guangdong_liberal/q15_fig1.png}
\end{problem}

\begin{answer}
\(\frac{\sqrt{21}}{2}\)
\end{answer}

\begin{solution}
建立直角坐标系，取 \(A=(0,0)\)，\(B=(0,\sqrt3)\)，\(C=(3,\sqrt3)\)，\(D=(3,0)\). 则 \(\overrightarrow{AC}=(3,\sqrt3)\)，且 \(E\) 是点 \(B\) 在直线 \(AC\) 上的投影.

设 \(E=t\overrightarrow{AC}\). 由 \(BE\perp AC\)，有
\[
\left(B-t\overrightarrow{AC}\right)\cdot\overrightarrow{AC}=0
\]
代入 \(B=(0,\sqrt3)\)，得
\[
t=\frac{B\cdot\overrightarrow{AC}}{|\overrightarrow{AC}|^2}=\frac{3}{12}=\frac14
\]
所以 \(E=\left(\frac34,\frac{\sqrt3}{4}\right)\). 于是
\[
ED^2=\left(3-\frac34\right)^2+\left(\frac{\sqrt3}{4}\right)^2=\frac{81}{16}+\frac3{16}=\frac{21}{4}
\]
因此 \(ED=\frac{\sqrt{21}}2\).
\end{solution}

\section{解答题}

\begin{problem}
已知函数 \(f(x) = \sqrt{2}\cos \left(x - \frac{\pi}{12}\right)\)，\(x \in \R\).
\begin{enumerate}
\item 求 \( f\left(\frac{\pi}{3}\right) \) 的值；
\item 若 \(\cos\theta=\frac{3}{5}\)，\(\theta\in\left(\frac{3\pi}{2},2\pi\right)\)，求 \( f\left(\theta-\frac{\pi}{6}\right) \).
\end{enumerate}
\end{problem}

\begin{solution}
\begin{enumerate}
\item
将 \(x=\frac{\pi}{3}\) 代入，得
\[
f\left(\frac{\pi}{3}\right)=\sqrt2\cos\left(\frac{\pi}{3}-\frac{\pi}{12}\right)=\sqrt2\cos\frac{\pi}{4}=1
\]

\item
因为 \(\theta\in\left(\frac{3\pi}{2},2\pi\right)\)，所以 \(\sin\theta<0\). 由 \(\cos\theta=\frac35\) 及 \(\sin^2\theta+\cos^2\theta=1\)，得 \(\sin\theta=-\frac45\). 于是
\[
\begin{aligned}
f\left(\theta-\frac{\pi}{6}\right)
&=\sqrt2\cos\left(\theta-\frac{\pi}{6}-\frac{\pi}{12}\right)\\
&=\sqrt2\cos\left(\theta-\frac{\pi}{4}\right)\\
&=\sqrt2\left(\cos\theta\cos\frac{\pi}{4}+\sin\theta\sin\frac{\pi}{4}\right)\\
&=\sqrt2\left(\frac{3}{5\sqrt2}-\frac{4}{5\sqrt2}\right)=-\frac15.
\end{aligned}
\]
\end{enumerate}
\end{solution}

\begin{problem}
从一批苹果中，随机抽取 \(50\) 个，其重量（单位：克）的频数分布表如下：

\begin{center}
\begin{tabular}{c|cccc}
分组（重量） & \( [80,85) \) & \( [85,90) \) & \( [90,95) \) & \( [95,100) \) \\
\hline
频数 & \(5\) & \(10\) & \(20\) & \(15\)
\end{tabular}
\end{center}

\begin{enumerate}
\item 根据频数分布表计算苹果的重量在 \([90,95)\) 的频率；
\item 用分层抽样的方法从重量在 \([80,85)\) 和 \([95,100)\) 的苹果中共抽取 4 个，其中重量在 \([80,85)\) 的有几个？
\item 在（2）中抽出的 4 个苹果中，任取 2 个，求重量在 \([80,85)\) 和 \([95,100)\) 中各有 1 个的概率.
\end{enumerate}
\end{problem}

\begin{solution}
\begin{enumerate}
\item
重量在 \([90,95)\) 中的频数为 \(20\)，样本总数为 \(50\)，所以该组的频率为
\[
\frac{20}{50}=0.4
\]

\item
重量在 \([80,85)\) 和 \([95,100)\) 中的苹果总数为 \(5+15=20\). 分层抽样时各层抽取数之比等于各层苹果数之比，因此重量在 \([80,85)\) 中应抽取
\[
4\times\frac{5}{20}=1
\]
个苹果.

\item
由（2）知抽出的 \(4\) 个苹果中，\([80,85)\) 中有 \(1\) 个，\([95,100)\) 中有 \(3\) 个. 任取 \(2\) 个的基本事件数为 \(\mathrm C_4^2=6\)，其中两层各取 \(1\) 个的事件数为 \(\mathrm C_1^1\mathrm C_3^1=3\). 所以所求概率为
\[
\frac{\mathrm C_1^1\mathrm C_3^1}{\mathrm C_4^2}=\frac{3}{6}=\frac12
\]
\end{enumerate}
\end{solution}

\begin{problem}
如图\circled{1}，在边长为 1 的等边三角形 \(ABC\) 中，\(D\)，\(E\) 分别是 \(AB\)，\(AC\) 上的点，\(AD = AE\)，\(F\) 是 \(BC\) 的中点，\(AF\) 与 \(DE\) 交于点 \(G\)，将 \(\triangle ABF\) 沿 \(AF\) 折起，得到如图\circled{2}所示的三棱锥 \(A - BCF\)，其中 \(BC = \frac{\sqrt{2}}{2}\).
\begin{enumerate}
\item 证明： \(DE \parallel\) 平面 \(BCF\)；
\item 证明： \(CF \perp\) 平面 \(ABF\)；
\item 当 \(AD = \frac{2}{3}\) 时，求三棱锥 \(F - DEG\) 的体积 \(V_{F - DEG}\).
\end{enumerate}

\bitmapfigure[max width=0.88\linewidth]{img/2013/guangdong_liberal/q18_fig1.png}
\end{problem}

\begin{solution}
\begin{enumerate}
\item
设 \(t=\frac{AD}{AB}=\frac{AE}{AC}\). 折叠后点 \(D\)、\(E\) 仍分别把 \(AB\)、\(AC\) 按相同比例分割，因此
\[
\overrightarrow{DE}=t\overrightarrow{BC}
\]
所以 \(DE\parallel BC\). 又 \(BC\subset\) 平面 \(BCF\)，且 \(DE\) 不在平面 \(BCF\) 内，故 \(DE\parallel\) 平面 \(BCF\).

\item
在原等边三角形中，\(F\) 是 \(BC\) 的中点，\(AF\) 是中线，也是高，所以 \(AF\perp BC\)，从而 \(AF\perp CF\). 折叠沿 \(AF\) 进行，三角形 \(AFC\) 未被折起，故折叠后 \(AF\)、\(CF\) 的位置不变，仍有 \(AF\perp CF\). 又 \(BF=CF=\frac12\)，且题设 \(BC=\frac{\sqrt2}{2}\)，从而
\[
BF^2+CF^2=\frac14+\frac14=\frac12=BC^2
\]
由勾股定理的逆定理，\(BF\perp CF\). 直线 \(AF\) 与 \(BF\) 在平面 \(ABF\) 内相交，且 \(CF\) 同时垂直于它们，所以 \(CF\perp\) 平面 \(ABF\).

\item
由（2），取 \(F\) 为原点，以 \(FA\)、\(FB\)、\(FC\) 所在方向为三个坐标轴，则
\(A=\left(\frac{\sqrt3}{2},0,0\right)\)，\(B=\left(0,\frac12,0\right)\)，\(C=\left(0,0,\frac12\right)\).
当 \(AD=\frac23\) 时，\(D\)、\(E\) 分别把 \(AB\)、\(AC\) 按 \(1:2\) 分割，所以
\(D=\frac13A+\frac23B=\left(\frac{\sqrt3}{6},\frac13,0\right)\)，\(E=\frac13A+\frac23C=\left(\frac{\sqrt3}{6},0,\frac13\right)\).
在原平面图中，由相似三角形得 \(\frac{AG}{AF}=\frac{AD}{AB}=\frac23\)，而 \(G\) 在折痕 \(AF\) 上不动，因此
\[
G=\left(\frac{\sqrt3}{6},0,0\right)
\]
故 \(\triangle DEG\) 是直角三角形，且 \(DG=EG=\frac13\)，其面积为 \(\frac1{18}\). 平面 \(DEG\) 的方程为 \(x=\frac{\sqrt3}{6}\)，所以点 \(F\) 到平面 \(DEG\) 的距离为 \(\frac{\sqrt3}{6}\). 于是
\[
V_{F-DEG}=\frac13\times\frac1{18}\times\frac{\sqrt3}{6}=\frac{\sqrt3}{324}
\]
\end{enumerate}
\end{solution}

\begin{problem}
设各项均为正数的数列 \(\{a_{n}\}\) 的前 \(n\) 项和为 \(S_{n}\)，满足 \(4S_{n} = a_{n+1}^{2} - 4n - 1\)，\(n \in \mathbf{N}^{*}\)，且 \(a_{2}\)，\(a_{5}\)，\(a_{14}\) 构成等比数列.
\begin{enumerate}
\item 证明： \(a_{2} = \sqrt{4a_{1} + 5}\)；
\item 求数列 \(\{a_n\}\) 的通项公式；
\item 证明：对一切正整数 \(n\)，有 \(\frac{1}{a_1 a_2} + \frac{1}{a_2 a_3} + \cdots + \frac{1}{a_n a_{n+1}} < \frac{1}{2}\).
\end{enumerate}
\end{problem}

\begin{solution}
\begin{enumerate}
\item
取 \(n=1\)，由 \(S_1=a_1\) 得
\[
4a_1=a_2^2-4-1=a_2^2-5
\]
因为各项均为正数，所以 \(a_2>0\)，从而 \(a_2=\sqrt{4a_1+5}\).

\item
对 \(n\geq2\)，将题设中 \(n\) 与 \(n-1\) 两式相减，得
\[
4a_n=a_{n+1}^2-a_n^2-4
\]
即 \(a_{n+1}^2=(a_n+2)^2\). 由各项为正，得 \(a_{n+1}=a_n+2\)，所以当 \(n\geq2\) 时
\[
a_n=a_2+2(n-2)
\]
设 \(a_2=b>0\)，则 \(a_5=b+6\)，\(a_{14}=b+24\). 由 \(a_2\)、\(a_5\)、\(a_{14}\) 成等比数列，得
\[
(b+6)^2=b(b+24)
\]
化简得 \(12b=36\)，所以 \(b=3\). 再由第（1）问的关系 \(9=4a_1+5\)，得 \(a_1=1\). 因此
\(a_n=2n-1\)，\((n\in\mathbf N^*)\).

\item
由 \(a_n=2n-1\)，有
\[
\frac1{a_na_{n+1}}=\frac1{(2n-1)(2n+1)}
=\frac12\left(\frac1{2n-1}-\frac1{2n+1}\right)
\]
所以
\[
\sum_{k=1}^{n}\frac1{a_ka_{k+1}}
=\frac12\left(1-\frac1{2n+1}\right)
=\frac{n}{2n+1}<\frac12
\]
\end{enumerate}
\end{solution}

\begin{problem}
已知抛物线 \(C\) 的顶点为原点，其焦点 \(F(0, c)\)（\(c > 0\)）到直线 \(l\)：\(x - y - 2 = 0\) 的距离为 \(\frac{3\sqrt{2}}{2}\). 设 \(P\) 为直线 \(l\) 上的点，过点 \(P\) 作抛物线 \(C\) 的两条切线 \(PA\)，\(PB\)，其中 \(A\)，\(B\) 为切点.
\begin{enumerate}
\item 求抛物线 \(C\) 的方程；
\item 当点 \(P(x_0, y_0)\) 为直线 \(l\) 上的定点时，求直线 \(AB\) 的方程；
\item 当点 \(P\) 在直线 \(l\) 上移动时，求 \(|AF| \cdot |BF|\) 的最小值.
\end{enumerate}
\end{problem}

\begin{solution}
\begin{enumerate}
\item
抛物线顶点为原点，焦点为 \(F(0,c)\)，且 \(c>0\)，故其标准方程为 \(x^2=4cy\). 由点 \(F\) 到直线 \(x-y-2=0\) 的距离，得
\[
\frac{|0-c-2|}{\sqrt2}=\frac{c+2}{\sqrt2}=\frac{3\sqrt2}{2}
\]
因此 \(c+2=3\)，即 \(c=1\)，所以抛物线方程为 \(x^2=4y\).

\item
令抛物线上参数为 \(t\) 的点为 \(Q(t)=(2t,t^2)\)，则该点处的切线方程为
\[
tx-y-t^2=0
\]
设两切点对应的参数为 \(t_1\)、\(t_2\). 由于两条切线都经过 \(P(x_0,y_0)\)，所以 \(t_1\)、\(t_2\) 是方程
\[
t^2-x_0t+y_0=0
\]
的两个根，从而 \(t_1+t_2=x_0\)，\(t_1t_2=y_0\). 过 \(Q(t_1)\)、\(Q(t_2)\) 的直线为
\[
(t_1+t_2)x-2y-2t_1t_2=0
\]
代入根的和与积，得直线 \(AB\) 的方程为
\[
x_0x-2y-2y_0=0
\]

\item
焦点为 \(F(0,1)\). 对参数为 \(t\) 的切点 \(Q(t)\)，有
\[
|Q(t)F|=\sqrt{(2t)^2+(t^2-1)^2}=t^2+1
\]
因此
\[
|AF|\cdot|BF|=(1+t_1^2)(1+t_2^2)
\]
记 \(p=t_1+t_2\)，\(q=t_1t_2\). 因为 \(P(x_0,y_0)\) 在直线 \(l\) 上，且 \(x_0-y_0-2=0\)，所以 \(p-q=2\)，即 \(q=p-2\). 于是
\[
\begin{aligned}
(1+t_1^2)(1+t_2^2)
&=1+(p^2-2q)+q^2\\
&=2p^2-6p+9\\
&=2\left(p-\frac32\right)^2+\frac92\geq\frac92.
\end{aligned}
\]
当 \(p=\frac32\) 时，取 \(t_1=\frac{3+\sqrt{17}}4\)、\(t_2=\frac{3-\sqrt{17}}4\)，则 \(t_1+t_2=\frac32\)、\(t_1t_2=-\frac12\). 相应的 \(P=(\frac32,-\frac12)\) 满足直线 \(l\) 的方程，且确有两条切线，等号可以取得. 因此最小值为 \(\frac92\).
\end{enumerate}
\end{solution}

\begin{problem}
设函数 \( f(x) = x^3 - kx^2 + x  \)（\(x \in \R\)）.
\begin{enumerate}
\item 当 \(k = 1\) 时，求函数 \(f(x)\) 的单调区间；
\item 当 \(k < 0\) 时，求函数 \(f(x)\) 在 \([k, -k]\) 上的最小值 \(m\) 和最大值 \(M\).
\end{enumerate}
\end{problem}

\begin{solution}
\begin{enumerate}
\item
当 \(k=1\) 时，\(f(x)=x^3-x^2+x\)，其导数为
\[
f'(x)=3x^2-2x+1=3\left(x-\frac13\right)^2+\frac23>0
\]
所以 \(f(x)\) 在 \(\R\) 上单调递增.

\item
当 \(k<0\) 时，\(f(k)=k\)，并且
\[
f(x)-f(k)=x^3-kx^2+x-k=(x-k)(x^2+1)
\]
当 \(x\in[k,-k]\) 时，\(x-k\geq0\)，且 \(x^2+1>0\)，故 \(f(x)\geq k\)，所以最小值 \(m=k\).

又 \(f(-k)=-2k^3-k\)，并且
\[
f(-k)-f(x)=(-k-x)(x^2-2kx+2k^2+1)
\]
当 \(x\in[k,-k]\) 时，\(-k-x\geq0\)，且
\[
x^2-2kx+2k^2+1=(x-k)^2+k^2+1>0
\]
故 \(f(x)\leq f(-k)=-2k^3-k\)，所以最大值 \(M=-2k^3-k\).
\end{enumerate}
\end{solution}
